\documentclass{article}
\usepackage[utf8]{inputenc}
\usepackage[german]{babel}
\usepackage{amsmath, amssymb, amsthm}
\usepackage{booktabs}
\usepackage{array}

\newtheorem{theorem}{Satz}
\newcolumntype{L}[1]{>{\raggedright\arraybackslash}p{#1}}

\title{Vellemans Beweis-Rezepte (Vollständige Vorher/Nachher-Ansicht)}
\author{Nach Daniel J. Velleman: \emph{How to Prove It}}
\date{\today}

\begin{document}

\maketitle

\section*{Vellemans Logische Rezepte (Scratch-Regeln)}

\subsection*{1. Die UND-Verknüpfung ($P \land Q$)}
\begin{center}
\small
\begin{tabular}{ L{0.45\textwidth} | L{0.45\textwidth} }
\toprule
\textbf{Wenn $P \land Q$ in den GIVENS steht} & \textbf{Wenn $P \land Q$ im GOAL steht} \\
\midrule
\textbf{Vorher:} \newline
1. $P \land Q$ \newline
\emph{Goal:} $R$ &
\textbf{Vorher:} \newline
1. [Voraussetzungen] \newline
\emph{Goal:} $P \land Q$ \\
\hline
\textbf{Nachher (Rezept angewendet):} \newline
1. $P$ \newline
2. $Q$ \newline
\emph{Goal:} $R$ \newline
(\emph{Man darf $P$ und $Q$ als separate Wahrheiten nutzen.}) &
\textbf{Nachher (Rezept angewendet):} \newline
Es entstehen zwei separate Beweis-Zweige: \newline
\textbf{Zweig 1:} Zeige Goal $P$. \newline
\textbf{Zweig 2:} Zeige Goal $Q$. \\
\bottomrule
\end{tabular}
\end{center}

\subsection*{2. Die IMPLIKATION ($P \rightarrow Q$)}
\begin{center}
\small
\begin{tabular}{ L{0.45\textwidth} | L{0.45\textwidth} }
\toprule
\textbf{Wenn $P \rightarrow Q$ in den GIVENS steht} & \textbf{Wenn $P \rightarrow Q$ im GOAL steht} \\
\midrule
\textbf{Vorher:} \newline
1. $P \rightarrow Q$ \newline
\emph{Goal:} $R$ &
\textbf{Vorher:} \newline
1. [Voraussetzungen] \newline
\emph{Goal:} $P \rightarrow Q$ \\
\hline
\textbf{Nachher (Modus Ponens):} \newline
Wenn du im Verlauf des Beweises irgendwann das Given $P$ erreichst, darfst du sofort hinzufügen: \newline
2. $Q$ &
\textbf{Nachher (Conditional Proof):} \newline
1. [Voraussetzungen] \newline
2. $P$ (\emph{Nimm die linke Seite an!}) \newline
\emph{Neues Goal:} $Q$ \\
\bottomrule
\end{tabular}
\end{center}

\subsection*{3. Der ALLQUANTOR ($\forall x \, P(x)$)}
\begin{center}
\small
\begin{tabular}{ L{0.45\textwidth} | L{0.45\textwidth} }
\toprule
\textbf{Wenn $\forall x \, P(x)$ in den GIVENS steht} & \textbf{Wenn $\forall x \, P(x)$ im GOAL steht} \\
\midrule
\textbf{Vorher:} \newline
1. $\forall x \, P(x)$ \newline
\emph{Goal:} $R$ &
\textbf{Vorher:} \newline
1. [Voraussetzungen] \newline
\emph{Goal:} $\forall x \, P(x)$ \\
\hline
\textbf{Nachher (Universal Instantiation):} \newline
Wenn im Beweis bereits ein konkretes Objekt (z.B. $c$) existiert, darfst du eintragen: \newline
2. $P(c)$ \newline
(\emph{Du ,,stopfst`` ein existierendes Objekt hinein.}) &
\textbf{Nachher (Universal Generalization):} \newline
1. [Voraussetzungen] \newline
2. Sei $c$ ein neues, \emph{beliebiges} Objekt. \newline
\emph{Neues Goal:} $P(c)$ \\
\bottomrule
\end{tabular}
\end{center}

\newpage

\subsection*{4. Die ODER-Verknüpfung ($P \lor Q$)}
\begin{center}
\small
\begin{tabular}{ L{0.45\textwidth} | L{0.45\textwidth} }
\toprule
\textbf{Wenn $P \lor Q$ in den GIVENS steht} & \textbf{Wenn $P \lor Q$ im GOAL steht} \\
\midrule
\textbf{Vorher:} \newline
1. $P \lor Q$ \newline
\emph{Goal:} $R$ &
\textbf{Vorher:} \newline
1. [Voraussetzungen] \newline
\emph{Goal:} $P \lor Q$ \\
\hline
\textbf{Nachher (Proof by Cases):} \newline
Der Beweis verzweigt sich in zwei Pfade. Du musst das Goal $R$ zweimal zeigen: \newline
\textbf{Fall 1:} Füge Given $P$ hinzu $\rightarrow$ Zeige $R$. \newline
\textbf{Fall 2:} Füge Given $Q$ hinzu $\rightarrow$ Zeige $R$. &
\textbf{Nachher (Rezept angewendet):} \newline
Wähle eine von zwei Strategien: \newline
1. Zeige direkt $P$ (oder direkt $Q$). \newline
2. Oder (häufiger): Nimm an, $P$ wäre falsch. \newline
$\rightarrow$ Given: $\neg P$ \newline
$\rightarrow$ \emph{Neues Goal:} $Q$ \\
\bottomrule
\end{tabular}
\end{center}

\subsection*{5. Der EXISTENZQUANTOR ($\exists x \, P(x)$)}
\begin{center}
\small
\begin{tabular}{ L{0.45\textwidth} | L{0.45\textwidth} }
\toprule
\textbf{Wenn $\exists x \, P(x)$ in den GIVENS steht} & \textbf{Wenn $\exists x \, P(x)$ im GOAL steht} \\
\midrule
\textbf{Vorher:} \newline
1. $\exists x \, P(x)$ \newline
\emph{Goal:} $R$ &
\textbf{Vorher:} \newline
1. [Voraussetzungen] \newline
\emph{Goal:} $\exists x \, P(x)$ \\
\hline
\textbf{Nachher (Existential Instantiation):} \newline
Führe einen Namen für das existierende Objekt ein (z. B. $x_0$). \newline
1. $x_0$ ist ein konkretes Objekt. \newline
2. $P(x_0)$ \newline
\emph{Goal:} $R$ \newline
(\emph{Wichtig: $x_0$ muss ein brandneuer Variablenname sein!}) &
\textbf{Nachher (Konstruktion):} \newline
Du musst aus deinen Givens einen konkreten Kandidaten (nennen wir ihn $c$) finden oder berechnen. \newline
1. [Voraussetzungen] \newline
\emph{Neues Goal:} $P(c)$ \newline
(\emph{Du musst beweisen, dass dein Kandidat die Eigenschaft erfüllt.}) \\
\bottomrule
\end{tabular}
\end{center}

\newpage

\section*{Anwendungsbeispiel für $\exists$ und $\lor$}

\begin{theorem}
Seien $A, B, C$ Mengen. Wenn $A \neq \emptyset$ und $A \subseteq B$, dann ist $A \cup C \neq \emptyset$.
\end{theorem}

\subsection*{Scratch Work}
Zur Erinnerung: $X \neq \emptyset$ bedeutet logisch übersetzt $\exists x \, (x \in X)$.

\begin{center}
\small
\begin{tabular}{ L{0.45\textwidth} | L{0.45\textwidth} }
\toprule
\textbf{Givens} & \textbf{Goals} \\
\midrule
1. $\exists x \, (x \in A)$ \newline
2. $\forall z \, (z \in A \rightarrow z \in B)$ &
\emph{Goal:} $\exists x \, (x \in A \lor x \in C)$ \\[1ex]
\hline
% Anwendung: Existenzquantor im Given abbauen
3. $x_0$ ist ein konkretes Objekt. \newline
4. $x_0 \in A$ \newline
(\emph{Nach Rezept: $\exists$ in Givens}) \newline
5. $\forall z \, (z \in A \rightarrow z \in B)$ &
\emph{Goal:} $\exists x \, (x \in A \lor x \in C)$ \\[1ex]
\hline
% Anwendung: Allquantor im Given nutzen
3. $x_0$ ist ein konkretes Objekt. \newline
4. $x_0 \in A$ \newline
5. $x_0 \in A \rightarrow x_0 \in B$ \quad (\emph{Instanziierung}) \newline
6. $x_0 \in B$ \quad (\emph{Modus Ponens}) &
\emph{Goal:} $\exists x \, (x \in A \lor x \in C)$ \\[1ex]
\hline
% Anwendung: Existenzquantor im Goal abbauen (Konstruktion)
% Wir wählen x_0 als unseren Zeugen!
3. $x_0$ ist ein konkretes Objekt. \newline
4. $x_0 \in A$ \newline
6. $x_0 \in B$ &
\emph{Neues Goal:} $x_0 \in A \lor x_0 \in C$ \newline
(\emph{Nach Rezept: $\exists$ im Goal mit Kandidate $x_0$}) \\[1ex]
\hline
% Anwendung: Oder im Goal erfüllen
3. $x_0$ ist ein konkretes Objekt. \newline
4. $x_0 \in A$ \newline
6. $x_0 \in B$ &
$x_0 \in A \lor x_0 \in C$ \quad \checkmark \newline
(\emph{Ziel erreicht, da Given 4 wahr ist!}) \\
\bottomrule
\end{tabular}
\end{center}

\subsection*{Formaler Beweis}
\begin{proof}
Angenommen, $A \neq \emptyset$ und $A \subseteq B$. Da $A \neq \emptyset$ gilt, existiert ein Element in $A$. Wir nennen dieses Element $x_0$ (es gilt also $x_0 \in A$). Da nach Voraussetzung $A \subseteq B$ gilt (also jedes Element aus $A$ auch in $B$ liegt), folgt aus $x_0 \in A$ sofort $x_0 \in B$.

Um zu zeigen, dass $A \cup C \neq \emptyset$ gilt, müssen wir zeigen, dass ein Element existiert, das in $A \cup C$ liegt. Wir wählen $x_0$ als unseren Kandidaten. Da bereits $x_0 \in A$ wahr ist, ist die Oder-Aussage $x_0 \in A \lor x_0 \in C$ ebenfalls wahr. Somit ist $x_0 \in A \cup C$, woraus folgt, dass $A \cup C$ nicht leer ist.
\end{proof}

\end{document}
